% Part:sets-functions-relations
% Chapter: size-of-sets
% Section: non-enumerability-alt

\documentclass[../../../include/open-logic-section]{subfiles}

\begin{document}

\olfileid{sfr}{siz}{nen-alt}

\olsection{Himpunan kang \printtoken{S}{nonenumerable}}

\begin{editorial}
  Bagean iki mbuktekake yen $\Bin^\omega$ lan
  $\Pow{\Nat}$ ora bisa didhaptar, nganggo definisi ing \olref[enm-alt]{sec},
  yaiku mbutuhake bijeksi karo~$\Nat$, dudu surjeksi saka
  $\PosInt$.
\end{editorial}

\begin{explain}
Himpunan $\Nat$ saka wilangan natural iku tanpa wates.
Himpunan iki uga cetha !!{enumerable}. Nanging kasunyatan kang
nggumunake yaiku anane himpunan kang \emph{!!{nonenumerable}},
yaiku himpunan kang ora !!{enumerable}
(delengen \olref[sfr][siz][enm-alt]{defn:enumerable}).

Iki bisa uga nggumunake. Ngandharake yen $A$ iku
!!{nonenumerable} ateges \emph{ora ana} !!{bijection} $f
\colon \Nat \to A$, yaiku ora ana fungsi kang metakake
!!{element}s saka~$\Nat$ kang cacahe tanpa wates menyang~$A$
lan nyakup kabeh~$A$. Dadi yen $A$ iku !!{nonenumerable},
!!{element}s saka~$A$ "luwih akeh" tinimbang wilangan natural.

Kanggo mbuktekake yen himpunan iku !!{nonenumerable}, kudu dituduhake
yen ora bisa ana !!{bijection} kang cocog. Cara kang paling becik
yaiku nuduhake yen saben upaya ngenumerasi !!{element}s saka~$A$
mesthi ora ngemot paling ora siji !!{element}. Iki nuduhake yen
ora ana fungsi $f\colon \Nat \to A$ kang !!{surjective}.
Cara umum kanggo mbuktekake iki yaiku nganggo \emph{metodhe diagonal} Cantor.
Saka dhaptar !!{element}s saka $A$, upamane $x_1$, $x_2$, \dots,
kita mbangun !!{element} saka~$A$ liyane kang, amarga cara anggone
dibangun, ora mungkin ana ing dhaptar mau.

Nanging kabeh iki luwih gampang dimangerteni nganggo tuladha.
Tuladha kapisan yaiku himpunan~$\Bin^\omega$ saka kabeh rerangken
tanpa wates saka $0$ lan $1$. (Lambang "$\Bin$" nuduhake biner,
lan bisa dianggep minangka himpunan rong anggota
$\{0,1\}$.)\oliflabeldef{sfr:card-arithmetic:card-opps:sec}{\footnote{Kanthi
luwih pas, kita kudu netepake yen $\Bin^\omega$ iku himpunan kabeh
rerangken-$\omega$ saka $0$ lan $1$, yaiku himpunan
$\funfromto{\omega}{\{0,1\}}$. Nanging tegese iki lagi bakal
cetha ing \olref[sfr][card-arithmetic][card-opps]{sec}.}{} Nanging
andharan kang rada longgar iki saiki kudune ora
nuwuhake bingung.}
\end{explain}

\begin{thm}
\ollabel{thm:nonenum-bin-omega}
$\Bin^\omega$~iku !!{nonenumerable}.
\end{thm}

\begin{proof}
Gatekna sembarang enumerasi subhimpunan saka $\Bin^\omega$.
Dadi ana dhaptar $s_{0}$, $s_{1}$, $s_{2}$, \dots{}
kanthi saben $s_n$ minangka rerangken tanpa wates saka $0$ lan~$1$.
Upamakna $s_n(m)$ iku digit kaping $m$ saka rerangken kaping $n$
ing dhaptar iki. Urutan indeks iki dibenerake manut larikan sabanjure
amarga sumber Inggris mbalikke indeks (OLSIZ-008). Saiki dhaptar iki bisa dianggep minangka larikan,
kanthi $s_n(m)$ ing baris kaping $n$ lan kolom kaping $m$:
\[
\begin{array}{c|c|c|c|c|c}
& 0 & 1 & 2 & 3 & \dots \\\hline
0 & \mathbf{s_{0}(0)} & s_{0}(1) & s_{0}(2) & s_0(3) & \dots \\\hline
1 & s_{1}(0)& \mathbf{s_{1}(1)} & s_1(2) & s_1(3) & \dots \\\hline
2 & s_{2}(0)& s_{2}(1) & \mathbf{s_2(2)} & s_2(3) & \dots \\\hline
3 & s_{3}(0)& s_{3}(1) & s_3(2) & \mathbf{s_3(3)} & \dots \\\hline
\vdots & \vdots & \vdots & \vdots & \vdots & \mathbf{\ddots}
\end{array}
\]
Saiki kita bakal mbangun rerangken tanpa wates $d$ saka $0$ lan $1$
kang ora ana ing dhaptar iki. Carane yaiku nemtokake saben anggotane,
yaiku nemtokake $d(n)$ kanggo kabeh $n \in \Nat$.
Kanthi intuisi, kita maca diagonal larikan ndhuwur saka ndhuwur mudhun
(mula diarani "metodhe diagonal"), banjur ngganti saben $1$ dadi $0$
lan saben $0$ dadi~$1$. Pandhuan bit iki dibenerake dadi rong pambalikan
komplementer manut rumus kasus amarga sumber Inggris mbaleni arah kapisan
(OLSIZ-009). Kanthi luwih abstrak, $d(n)$ ditetepake dadi
$0$ utawa $1$ miturut apa !!{element} kaping $n$ ing diagonal,
yaiku $s_n(n)$, iku $1$ utawa $0$, yaiku:
\[
d(n) =
\begin{cases}
1 & \text{yen $s_{n}(n) = 0$}\\
0 & \text{yen $s_{n}(n) = 1$}
\end{cases}
\]
Cetha yen $d \in \Bin^\omega$, amarga iki rerangken tanpa wates
saka $0$ lan $1$. Nanging $d$ dibangun saengga $d(n) \neq s_n(n)$
kanggo saben $n \in \Nat$. Tegese, $d$ beda karo $s_n$
ing anggota kaping $n$. Mula $d \neq s_n$ kanggo saben $n\in \Nat$.
Dadi $d$ ora bisa ana ing dhaptar
$s_0$, $s_1$, $s_2$,
\dots

Kita wis nuduhake yen sembarang enumerasi sawijining subhimpunan
saka $\Bin^\omega$ mesthi ora ngemot sawijining !!{element} saka $\Bin^\omega$.
Dadi ora ana enumerasi himpunan $\Bin^\omega$, yaiku $\Bin^\omega$
iku !!{nonenumerable}.
\end{proof}

\begin{explain}
Cara pambuktian iki diarani "diagonalisasi" amarga nggunakake diagonal
larikan kanggo netepake~$d$. Nanging diagonalisasi ora kudu
nggunakake larikan. Kita bisa nuduhake yen sawijining himpunan
iku !!{nonenumerable} nganggo gagasan kang padha, sanajan ora ana larikan
utawa diagonal kang nyata. Asil ing ngisor nuduhake carane.
\end{explain}

\begin{thm}
\ollabel{thm:nonenum-pownat}
$\Pow{\Nat}$ ora !!{enumerable}.
\end{thm}

\begin{proof}
Kita nggunakake cara kang padha, yaiku nuduhake yen saben dhaptar
subhimpunan saka~$\Nat$ mesthi ora ngemot sawijining subhimpunan $\Nat$.
Upamakna ana dhaptar $N_0, N_1, N_2, \ldots$ saka subhimpunan $\Nat$.
Kita netepake himpunan $D$ kanthi: $n \in D$ yen lan mung yen $n \notin N_{n}$:
\[
D = \Setabs{n \in \Nat}{n \notin N_n}
\]
Cetha yen $D\subseteq \Nat$. Nanging $D$ ora bisa ana ing dhaptar.
Miturut cara pambangune, $n \in D$ yen lan mung yen $n\notin N_n$,
saengga $D \neq N_n$ kanggo saben $n \in \Nat$.
\end{proof}

\begin{explain}
Bukti sadurunge ora nyebut diagonal. Nanging kita bisa nganggep
ana diagonal yen nggambarake kaya mangkene: bayangna himpunan
$N_0$, $N_1$, \dots, ditulis ing larikan. Kita nulis $N_n$ ing
baris kaping $n$ kanthi nulis $m$ ing kolom kaping $m$
yen lan mung yen $m \in N_n$. Upamane, papat himpunan kapisan
ing dhaptar yaiku $\{0,1,2,\dots\}$, $\{1, 3, 5, \dots\}$, $\{0,1,4\}$,
lan $\{2,3,4,\dots\}$; mula wiwitan larikan:
\[
\begin{array}{r@{}rrrrrrr}
  N_0 = \{ & \mathbf{0}, & 1, & 2, & & & & \dots\}\\
  N_1 = \{ &  & \mathbf{1}, &  & 3, &  & 5, & \dots\}\\
  N_2 = \{ & 0, & 1, &  &  & 4\phantom{,} &  & \}\\
  N_3 = \{ &  &  & 2, & \mathbf{3}, & 4, & & \dots\}\\
  &\vdots & & & & & & \ddots\phantom{\}}
  \end{array}
\]
Dadi $D$ iku himpunan kang diolehake kanthi ngetutake diagonal mudhun,
kanthi netepake $n \in D$ yen lan mung yen $n$ \emph{ora} ana ing diagonal.
Ing tuladha ndhuwur, kita ora nglebokake $0$ lan $1$, nglebokake~$2$,
ora nglebokake~$3$, lan sateruse.
\end{explain}

\begin{prob}
Tuduhna yen himpunan kabeh fungsi $f \colon \Nat \to \Nat$
iku !!{nonenumerable} nganggo argumen diagonal kang eksplisit. Tegese,
tuduhna yen $f_1$, $f_2$, \dots, iku dhaptar fungsi lan saben $f_i\colon
\Nat \to \Nat$, mula ana $g \colon \Nat \to
\Nat$ kang ora ana ing dhaptar iki.
\end{prob}

\end{document}
